\section{分析}
\label{sec:analysis}

\subsection{为何跨模型移植有效？}
\label{sec:analysis:crossmodel}

Llama-2-chat 与 Llama-2-base 共享预训练初始化；RLHF 沿现有残差流基底移动激活而不旋转它 \cite{arditi2024refusal}。先前工作已实证确认：将 chat 模型的拒绝方向移植至同源 base LLM（相同预训练，无 RLHF）可将拒绝率从接近零提升至 88\% 以上 \cite{baselmsrefuse2024}。对于我们的场景，剩余的问题是这种共享几何是否能在 \textsc{OpenVLA} 额外 27 个 epoch 的行为克隆微调后得以保留。

为验证这一点，我们在 $1\mathrm{k}$ 自然语言探针上计算 Llama-2-chat 与 \textsc{OpenVLA} 主干之间第 $L$ 层隐层状态的平均余弦相似度。第 8--20 层的高相似度将表明基底实质上被保留，方向在两个模型间仍可互译。我们还检验从 Llama-2-chat 提取的 $r_L$ 在 \textsc{OpenVLA} 的有害动作位置上获得的零运动质量是否显著高于等范数随机单位向量（\S\ref{sec:exp:ablation} 方向来源对照）。较大的 $\Delta_{\text{rand}}$ 差值是几何保留的功能性印证：拒绝方向的特定几何性质——而非仅仅是其幅度——是有效成分。这些数值将在相机就绪版本中报告。

\subsection{RDT{}+ 下安全层探针的恢复}
\label{sec:analysis:probe-recovery}

\S\ref{sec:method:diagnostic} 中的诊断确立了拒绝方向在 \textsc{OpenVLA} 的动作 token 位置崩溃（投影 AUC 较低）。我们还在施加 \method{}+ 前后分别运行 \citet{securitytensors2025} 的安全层探针——该探针逐层测量良性-良性与良性-有害隐层状态对的余弦相似度差值。若动作 token 位置的探针分离度从接近零恢复至接近文本 token 位置的水平，则我们获得了对干预的机制层面而非行为层面的确认：该方向在恢复一个几何上缺失的信号，而不仅仅是任意破坏动作生成。该分析与解耦诊断结果一并报告。

\subsection{OpenVLA-OFT 作为方案特有负对照}
\label{sec:analysis:oft}

\textsc{OpenVLA-OFT} \cite{openvlaoft2025} 将最后 256 token 的离散化替换为连续 L1 回归头。按照我们的假设，\method{} 和 \method{}+ 是方案特有的：不存在 id $\geq 31744$ 的动作 token 可供定向。因此，我们预期将 \method{}+（相同 $r_L$、层、$\alpha$）施加于 \textsc{OpenVLA-OFT} 所产生的 HRR 变化与无防御基线无异。若此预测得到确认，则验证了 \method{} 的收益来自动作词汇交互，而非来自一般隐层状态扰动——对 \S\ref{sec:exp:ablation} 中方向来源对照实验的互补证据。

\subsection{拒绝在动作空间中意味着什么？}
\label{sec:analysis:action-space}

当 \method{}+ 在 $\alpha = 1$ 时触发，\textsc{OpenVLA} \emph{希望}输出什么来替代有害动作？我们探测 \method{}+ 关闭和开启时，第一个生成动作 token 在 256 个动作 bin 上的 softmax 分布。可能的定性结果有三种：（i）概率质量向 bin $128$（零增量，原地不动）移动；（ii）概率质量广泛重分布，表明 \method{}+ 仅破坏了有害动作而未引导出连贯的安全替代；（iii）概率质量向另一个仍然危险的区域移动。结果（i）支持将 \method{} 的拒绝解读为\emph{保持位置}——一种物理上安全的默认行为——而非注入随机运动。对夹爪自由度施加相同探测，可表征 \method{}+ 是否倾向于维持当前夹爪状态。我们还将该分布与方向来源变体（拒绝方向、随机方向、取反方向）进行对比，确认拒绝方向的特定几何性质是产生连贯停止行为的原因。

\begin{figure}[t]
\centering
\framebox[\columnwidth]{%
  \begin{minipage}{0.9\columnwidth}\centering\vspace{14pt}
  \emph{占位符图。} 有害提示下的动作 bin softmax 分布，对比无防御 / \method{}（仅动作）/ \method{}+（文本+动作）/ 随机方向对照。\vspace{14pt}
  \end{minipage}}
\caption{动作空间中的拒绝：\method{}+ 下第一个动作 token 分布向零运动 bin 的预测偏移，以及与随机方向对照的对比。}
\label{fig:action-space-refusal-zh}
\end{figure}

\subsection{自适应攻击者下的攻击成功率上界}
\label{sec:analysis:adaptive-bound}

在 \S\ref{sec:exp:adaptive} 的 PGD 正交化攻击下，攻击者每层最多可将 $h_L \cdot r_L$ 减小 $\epsilon \cdot \|J_\text{vis}\|_\text{op}$，其中 $J_\text{vis}$ 是在扰动补丁处视觉到 LLM 前向传播的 Jacobian，$\epsilon$ 为像素空间预算。实践中，该上界转化为有界的 HRR 下降，我们将连同自适应压力下测得的残余防御收益一并报告。

\subsection{观察到的局限性}
\label{sec:analysis:limitations}

我们预期两种削弱 \method{} 效果的情况：（i）当有害指令语义较轻微时（例如"把软玩具扔向婴儿"），拒绝方向的激活较弱，产生部分而非零运动动作；（ii）当视觉场景模糊时（例如塑料玩具刀），\method{}+ 可能误判并过度拒绝。两种情况在一定程度上由内容自适应 $\alpha$ 头（\S\ref{sec:method:adaptive}）缓解，但均未得到完全解决，这为局限性部分提供了动机。
